\chapter{Modele propose}

\section{Vue generale}
Le modele propose repose sur une architecture modulaire composee de modules logiciels specialises, appeles agents fonctionnels dans ce memoire. Chaque agent possede un role limite et produit une sortie exploitable par l'agent suivant. Cette separation reduit le couplage entre les etapes et facilite l'audit du traitement. Le terme agent est donc utilise au sens architectural et non au sens d'un agent cognitif autonome capable de negocier des objectifs ou de prendre des decisions independantes.

\begin{figure}[H]
  \centering
  \includegraphics[height=0.78\textheight,keepaspectratio]{fig_architecture_logminer_mermaid.png}
  \caption{Vue generale de l'architecture logique de \logminer{} sous forme de diagramme Mermaid.}
  \label{fig:architecture}
\end{figure}

La chaine generale est la suivante : collecte des journaux, parsing, normalisation, routage vers une famille, detection par modele compatible, correlation des anomalies candidates, affichage dans le dashboard et journalisation d'audit.

\section{Agents fonctionnels}
Le tableau \ref{tab:agents} resume les agents principaux.

\begin{table}[H]
\centering
\caption{Agents fonctionnels de \logminer}
\label{tab:agents}
\begin{tabularx}{\textwidth}{p{3cm}X}
\toprule
\textbf{Agent} & \textbf{Responsabilite} \\
\midrule
Collecteur & Recupere les journaux locaux, exports ou fichiers de datasets. \\
Parseur & Detecte le format et extrait les champs pertinents. \\
Normaliseur & Convertit les champs vers un schema commun. \\
Routeur & Identifie la famille de journaux et choisit le modele compatible. \\
Detecteur & Applique le modele de detection ou de classification. \\
Correlateur & Regroupe les anomalies candidates en incidents lisibles. \\
Dashboard & Affiche les resultats, filtres, graphes, ressources et decisions. \\
Audit & Trace les executions, decisions analyste et messages agents. \\
\bottomrule
\end{tabularx}
\end{table}

\section{Schema commun}
Le schema commun retient les champs necessaires a l'analyse transversale : horodatage, severite, hote, source, categorie, utilisateur, message, indicateur d'anomalie, score et metadonnees. Ce schema n'elimine pas la richesse des sources originales, mais fournit un noyau stable pour la detection et la visualisation.

\subsection{Principes de normalisation}
La normalisation suit trois principes. Le premier est la conservation : une information inconnue n'est pas supprimee, mais placee dans les metadonnees ou dans le message original. Le deuxieme est la comparabilite : les champs communs doivent permettre de filtrer plusieurs sources avec les memes criteres. Le troisieme est la tracabilite : chaque evenement normalise doit rester rattachable a sa source et a sa ligne d'origine lorsque cette information existe.

\subsection{Familles de journaux}
Les familles de journaux representent des contextes d'analyse distincts. Un journal Windows ne se traite pas comme un flux CICIDS2017, et une alerte Wazuh ne porte pas le meme sens qu'une ligne Apache brute. Le modele propose separe donc les familles pour adapter les extracteurs, les variables numeriques et les detecteurs.

\section{Routage multi-modeles}
Le routage est un element central de l'architecture. Plutot que d'appliquer un modele unique a tous les journaux, \logminer{} selectionne une famille de traitement. Les familles considerees incluent notamment Windows, Wazuh, Linux/auth, reseau, HDFS, BGL et fallback.

Le routage s'appuie sur des indices explicites : noms de colonnes, structure des messages, champs caracteristiques, type de fichier et distribution des valeurs. Lorsqu'aucune famille n'est suffisamment probable, le systeme utilise une route de fallback afin de conserver les donnees sans les perdre silencieusement.

\begin{figure}[H]
  \centering
  \includegraphics[width=0.9\textwidth]{fig_family_routing_ablation.png}
  \caption{Ablation operationnelle du routage par famille.}
  \label{fig:routing}
\end{figure}

\section{Modeles de detection}
Deux categories de modeles sont utilisees. Les modeles supervises RandomForest sont entraines sur des datasets labelises tels que Linux/auth, CICIDS2017 et CIC-DDoS2019. Les modeles non supervises, notamment Isolation Forest, sont utilises lorsque les labels fiables ne sont pas disponibles. Dans ce cas, les sorties sont des anomalies candidates.

\subsection{Decision et interpretation}
La sortie d'un modele est convertie en objet d'analyse comprenant au minimum la famille, le score, la decision, la severite estimee et les champs ayant motive l'affichage dans le dashboard. Cette representation evite de presenter un score isole. Elle favorise une lecture par incident, dans laquelle l'analyste peut comparer le message original, le contexte temporel et l'historique des decisions.

\begin{figure}[H]
  \centering
  \includegraphics[width=0.9\textwidth]{fig_model_portfolio_scale.png}
  \caption{Portefeuille de modeles et volumes associes.}
  \label{fig:portfolio}
\end{figure}

\section{Correlation et incidents}
Une anomalie candidate isolee n'est pas forcement une intrusion. Le correlateur regroupe donc les signaux selon la severite, la proximite temporelle, la source, la categorie et les messages. L'objectif est de produire des incidents plus exploitables qu'une liste brute de scores.

\section{Bus, audit et extensions}
La version locale s'appuie sur des fichiers JSONL et CSV reproductibles. La V2 FastAPI expose des endpoints de traitement et un dashboard web. Redis Streams est ajoute comme bus operationnel local pour l'orchestration evenementielle et la repartition de taches entre workers. MQTT reste une extension de communication legere. Ces extensions ne sont pas presentees comme une preuve de scalabilite industrielle, mais comme une validation locale de l'architecture et une trajectoire d'evolution.

\section{Positionnement du modele}
Le modele propose est un complement aux outils SOC/SIEM existants. Il ne remplace pas Wazuh, OSSEC ou fail2ban. Il ajoute une couche analytique permettant de tester des detecteurs legers, de router les familles de journaux et de fournir un dashboard d'aide a l'analyse.

\section{Contraintes de conception}
La conception est guidee par quatre contraintes. La premiere est la reproductibilite : les resultats doivent pouvoir etre regeneres par scripts. La deuxieme est la legerete : le prototype doit fonctionner localement sans infrastructure SOC complete. La troisieme est l'auditabilite : les decisions et les messages agents doivent etre consultables. La quatrieme est la prudence scientifique : le systeme distingue les anomalies candidates, les alertes issues d'outils externes et les intrusions confirmees.

\section{Flot de traitement detaille}
Le flot de traitement commence par l'identification de la source. Le systeme peut recevoir un fichier local, un export de plateforme, un dataset public ou une sortie produite par un script de collecte. La premiere operation consiste a associer cette entree a un type probable. Cette etape ne depend pas d'un seul indice ; elle peut utiliser l'extension du fichier, les noms de colonnes, la structure des lignes, la presence de champs caracteristiques ou la distribution de certaines valeurs.

La deuxieme operation est le parsing. Le parseur transforme les lignes brutes en objets intermediaires. Lorsque le format est reconnu, les champs importants sont extraits. Lorsque le format est partiellement reconnu, les champs disponibles sont conserves et le message brut reste accessible. Lorsque le format est inconnu ou corrompu, le systeme conserve l'entree dans une route de fallback. Cette decision evite de supprimer des evenements qui pourraient signaler un probleme de format, une tentative d'evasion ou une erreur de collecte.

La troisieme operation est la normalisation. Les objets issus du parsing sont convertis dans un schema commun. Ce schema sert ensuite au routage et a l'affichage. La quatrieme operation est la detection. Le routeur determine la famille de logs, puis le detecteur applique le modele disponible. La cinquieme operation est la correlation. Elle regroupe les signaux proches par temps, source, severite ou categorie afin de produire des incidents plus comprehensibles qu'une liste brute d'anomalies.

\section{Modele logique des interactions}
Les agents cooperent selon un modele de messages. Chaque message transporte un identifiant, une source, un horodatage, une famille probable, un contenu normalise et des metadonnees. Cette logique permet de separer la production, la transformation et la consommation des informations. Dans la version locale, cette cooperation peut etre materialisee par des appels de fonctions, des fichiers, des endpoints FastAPI ou Redis Streams. Dans une version distribuee plus avancee, cette logique pourrait etre etendue a plusieurs machines avec files persistantes securisees.

La logique de message presente un avantage important : elle rend l'architecture observable. Lorsqu'un resultat semble incoherent, il devient possible de remonter a l'etape de parsing, au choix du routeur ou au modele utilise. Cette observabilite interne est aussi importante que la detection elle-meme, car un systeme de securite doit pouvoir expliquer ses propres decisions.

\section{Gestion des erreurs et cas limites}
Le modele prevoit plusieurs cas limites. Si un fichier ne contient pas les colonnes attendues, le systeme ne doit pas s'arreter brutalement. Si une ligne est corrompue, elle doit etre isolee et conservee. Si aucun modele n'est disponible pour une famille, le systeme doit produire un statut explicite plutot qu'un resultat trompeur. Si le dashboard ne peut pas charger une ressource, l'erreur doit etre visible dans l'audit.

Ces choix peuvent sembler secondaires, mais ils sont essentiels pour une application de cybersecurite. Les donnees reelles sont rarement propres. Les logs peuvent contenir des caracteres inattendus, des timestamps manquants, des champs vides ou des volumes tres variables. La robustesse consiste donc a degrader proprement le traitement, pas a supposer que toutes les entrees seront parfaites.

\section{Modele de donnees conceptuel}
Le modele de donnees conceptuel distingue quatre niveaux. Le premier niveau est l'evenement brut. Il correspond a la ligne, a l'enregistrement ou a l'objet recu depuis une source. A ce niveau, le systeme ne suppose pas encore que les champs sont fiables ou complets. Le deuxieme niveau est l'evenement parse. Il contient les champs extraits par le parseur, mais il reste proche de la source d'origine. Le troisieme niveau est l'evenement normalise. Il respecte le schema commun et peut etre manipule par les agents de detection, de correlation et de visualisation. Le quatrieme niveau est l'incident. Il regroupe une ou plusieurs anomalies candidates avec un contexte lisible.

Cette separation evite de melanger plusieurs notions. Une ligne de log n'est pas automatiquement une anomalie. Une anomalie candidate n'est pas automatiquement un incident. Un incident n'est pas automatiquement une intrusion confirmee. Chaque niveau ajoute une interpretation supplementaire. Cette progression est utile pour l'audit : lorsqu'une decision est contestee, il est possible de revenir au niveau precedent et de verifier si le probleme vient du parsing, du routage, du modele ou de la correlation.

Le schema commun joue un role d'interface entre les niveaux. Il contient les champs indispensables a l'analyse transversale : horodatage, source, hote, famille, categorie, severite, message, score, decision et metadonnees. Les champs absents ne doivent pas bloquer tout le traitement. Ils peuvent rester vides ou etre remplaces par une valeur explicite indiquant l'indisponibilite. Ce choix rend le systeme tolerant aux formats incomplets.

\section{Contrat de message entre modules}
L'architecture modulaire suppose que les composants echangent des informations sous forme de messages. Un message minimal contient un identifiant, un timestamp de production, une source, un type d'evenement, un payload et des metadonnees. L'identifiant permet de tracer un evenement dans la chaine. Le timestamp permet de reconstruire l'ordre des traitements. La source indique l'origine. Le payload transporte les champs utiles. Les metadonnees documentent les decisions intermediaires.

Ce contrat de message permet de separer les responsabilites. Le collecteur produit des messages bruts ou semi-structures. Le parseur enrichit le payload avec les champs extraits. Le normaliseur aligne les noms de colonnes. Le routeur ajoute la famille probable et le modele selectionne. Le detecteur ajoute le score et la decision. Le correlateur ajoute un identifiant d'incident. L'audit conserve les traces de ces transformations.

Dans la version locale, ces messages peuvent etre materialises par des appels de fonctions, des fichiers JSONL ou des sorties CSV. Dans une version plus distribuee, le meme contrat pourrait circuler dans Redis Streams, MQTT ou une autre file de messages. L'interet de cette abstraction est de preparer l'evolution sans imposer immediatement une infrastructure lourde.

\section{Strategie de correlation}
La correlation transforme des anomalies candidates isolees en incidents. Elle repose sur une hypothese simple : un evenement rare devient plus significatif lorsqu'il s'inscrit dans un contexte. Ce contexte peut etre temporel, spatial, semantique ou operationnel. Un echec de connexion isole peut etre banal. Plusieurs echecs rapproches sur le meme utilisateur ou la meme source peuvent devenir suspects. Une alerte Wazuh de severite moyenne peut devenir prioritaire si elle coincide avec une anomalie detectee par un modele non supervise.

La strategie retenue reste volontairement simple afin de conserver la lisibilite du prototype. Les anomalies sont regroupees selon la proximite temporelle, la source, la famille, la severite et certaines similarites de message. Cette correlation n'est pas encore un moteur de regles SOC complet. Elle constitue une premiere couche de regroupement qui reduit la fragmentation des resultats.

Une extension future consisterait a introduire des graphes d'evenements. Les noeuds representeraient les utilisateurs, hotes, adresses IP, processus, fichiers ou regles. Les aretes representeraient les relations observees. Une telle approche permettrait de reconstruire des chemins d'attaque ou des campagnes. Elle demanderait toutefois des donnees plus riches et une evaluation specifique.

\section{Modele de decision analyste}
Le modele propose integre explicitement la decision humaine. Le dashboard permet de valider, rejeter ou reclasser une anomalie candidate. Cette decision n'est pas seulement une annotation visuelle. Elle peut devenir un signal de retour. Dans une version future, les validations pourraient alimenter un apprentissage actif, ajuster les seuils ou enrichir les regles de correlation.

La decision analyste suit une logique prudente. L'analyste consulte d'abord le contexte : source, message, score, severite, famille, horodatage et incident associe. Il compare ensuite l'evenement avec les autres signaux de la meme periode. Il peut enfin choisir une action : confirmer l'interet, rejeter le signal, le reclasser ou le laisser en observation. L'audit conserve cette decision afin de permettre une relecture.

Ce modele reconnait que la detection d'anomalies est rarement totalement automatique. Dans un environnement reel, les alertes doivent etre interpretees. La valeur du systeme depend donc de sa capacite a fournir les bonnes informations au bon moment. L'agent dashboard est ainsi un composant du modele scientifique, pas seulement une interface graphique.

\section{Scenarios d'utilisation}
Plusieurs scenarios d'utilisation structurent le modele. Le premier est le scenario de traitement hors ligne. Un analyste fournit un fichier ou un dossier de logs. Le pipeline detecte les formats, normalise les evenements, applique les modeles et produit des CSV et tableaux. Ce scenario est utile pour l'analyse post-incident, la preparation de datasets et la redaction scientifique.

Le deuxieme scenario est le scenario de traitement local par micro-lots. Le serveur FastAPI expose un endpoint de decouverte et d'analyse. Le dashboard interroge les resultats et affiche les anomalies candidates. Ce scenario valide une boucle interactive periodique, mais reste limite a une machine ou a un environnement local. Il ne constitue pas une preuve de temps reel strict au sens d'un SOC a tres haut debit.

Le troisieme scenario est le scenario d'extension evenementielle. Les agents peuvent publier ou consommer des messages via Redis Streams ou MQTT. Redis Streams est utilise dans les experiences pour valider une repartition locale multi-processus, avec reprise d'une tache non acquittee. MQTT prepare une communication legere entre composants. Dans le memoire, ce scenario est presente comme une preuve locale d'architecture distribuee, non comme une preuve industrielle multi-machine.

Le quatrieme scenario est le scenario pedagogique. Le prototype sert a comprendre le cycle complet d'analyse des journaux : collecte, parsing, normalisation, detection, interpretation et audit. Ce scenario est important dans un cadre universitaire car il relie les notions theoriques de detection d'anomalies a un systeme concret.
